% !TEX TS-program = pdflatex
% !TeX program = pdflatex
% !TEX encoding = UTF-8
% !TeX spellcheck = en_US

\documentclass[xcolor=table]{beamer}

\usepackage{../extra/beamer/karimml}

\input{options2}

\subtitle[AE \& GANs]{Unsupervised deep learning}  

\changegraphpath{../img/unsupervised/}



\begin{document}
	
	\begin{frame}
		\frametitle{\inserttitle}
		\framesubtitle{\insertshortsubtitle: Motivation}
		
		
	\end{frame}
	
	\begin{frame}
		\frametitle{\inserttitle}
		\framesubtitle{\insertshortsubtitle: Some humor}
		
		\begin{center}
			\hgraphpage[0.9\textwidth]{humor/humor-unsupervised.jpg}
		\end{center}
		
	\end{frame}
	
	\begin{frame}
		\frametitle{\inserttitle}
		\framesubtitle{\insertshortsubtitle: Plan}
		
		\begin{multicols}{2}
			%	\small
			\tableofcontents
		\end{multicols}
	\end{frame}
	
	\section{Autoencoders}
	
	\begin{frame}
		\frametitle{\insertshortsubtitle}
		\framesubtitle{\insertsection}
		
		\begin{itemize}
			\item It is a multi-layer neural network.
			\item It learns a compressed latent representation of the data.
			\item It is characterized by \cite{2016-keras}:
			\begin{itemize}
				\item Data-driven: unlike compression algorithms like JPEG.
				\item Lossy reconstruction: the reconstructed output is not exactly identical to the input.
				\item Unsupervised learning.
			\end{itemize}
			\item Unlike specialized codecs (e.g., JPEG), autoencoders are generally not optimized for practical data compression.
		\end{itemize}
		
	\end{frame}
	
	\begin{frame}
		\frametitle{\insertshortsubtitle}
		\framesubtitle{\insertsection: Architecture}
		
		\hgraphpage[\textwidth]{AE.pdf}
		
	\end{frame}
	
	\subsection{Denoising autoencoder}
	
	\begin{frame}
		\frametitle{\insertshortsubtitle: \insertsection}
		\framesubtitle{\insertsubsection}
		
		\hgraphpage[\textwidth]{AE-noise.pdf}
		
	\end{frame}
	
	\begin{frame}
		\frametitle{\insertshortsubtitle: \insertsection}
		\framesubtitle{\insertsubsection: Some application}
		
		\begin{itemize}
			\item Improving the quality of speech (audio) \cite{2013-lu}.
			\item Cleaning dirty documents (refer to this competition: \url{https://www.kaggle.com/c/denoising-dirty-documents}).
			\item Retrieving historical documents \cite{2019-neji}.
			\item Completing hidden parts of the face \cite{2017-li-al}.
		\end{itemize}
		
	\end{frame}
	
	\subsection{Sparse autoencoder}
	
	\begin{frame}
		\frametitle{\insertshortsubtitle: \insertsection}
		\framesubtitle{\insertsubsection: Architecture}
		
		\hgraphpage[\textwidth]{AE-sparse.pdf}
		
	\end{frame}
	
	\begin{frame}
		\frametitle{\insertshortsubtitle: \insertsection}
		\framesubtitle{\insertsubsection: Description}
		
		\begin{itemize}
			\item Often use hidden layers equal to or larger than the input size.
			
			\item They automatically learn useful latent features for downstream tasks such as classification.
			
			\item Without constraints, the network may learn a trivial identity mapping:
			\[
			x_i \rightarrow h_i \rightarrow x_i
			\]
			
			\item Given an encoder $f$ and a decoder $g$, a sparsity penalty $\Omega$ is added to the hidden representation:
			\[
			\min \; J(x, g(f(x))) + \Omega(f(x))
			\]
			
			\item $\Omega$ is commonly based on: L1 sparsity, and Kullback--Leibler divergence.
			
			\item Classical sparse autoencoders usually use a single hidden layer.
			
			\item Deep sparse autoencoders can use multiple hidden layers to learn hierarchical features.
			
		\end{itemize}
		
	\end{frame}
	
	\begin{frame}
		\frametitle{\insertshortsubtitle: \insertsection}
		\framesubtitle{\insertsubsection: KL sparsity regularization}
		
		\begin{itemize}
			\item Sparse autoencoders encourage hidden neurons to activate rarely.
			
			\item Let:
			\[
			\hat{\rho}_j
			=
			\frac{1}{m}\sum_{i=1}^{m} h_j(x^{(i)})
			\]
			
			be the average activation of hidden neuron $j$.
			
			\item A small target activation $\rho$ is imposed:
			\[
			\rho \ll 1
			\]
			
			\item KL divergence penalizes deviations from the target sparsity:
			
			\[
			KL(\rho \,||\, \hat{\rho}_j)
			=
			\rho \log \frac{\rho}{\hat{\rho}_j}
			+
			(1-\rho)\log\frac{1-\rho}{1-\hat{\rho}_j}
			\]
			
			\item The final objective becomes:
			
			\[
			J_{sparse}
			=
			J_{reconstruction}
			+
			\beta \sum_j KL(\rho \,||\, \hat{\rho}_j)
			\]
			
		\end{itemize}
		
	\end{frame}
	
	\subsection{Variational autoencoder}
	
	\begin{frame}
		\frametitle{\insertshortsubtitle: \insertsection}
		\framesubtitle{\insertsubsection: Motivation}
		
		\begin{itemize}
			\item Standard autoencoders learn deterministic latent representations.
			
			\item Similar inputs are often separated into isolated latent clusters.
			
			\item Random points sampled between clusters may not correspond to valid data.
			
			\item This makes content generation difficult.
			
			\item \textbf{Idea of VAE:}
			\begin{itemize}
				\item learn a smooth probabilistic latent space,
				\item allow sampling and interpolation,
				\item generate new realistic data.
			\end{itemize}
		\end{itemize}
		
	\end{frame}
	
	
	\begin{frame}
		\frametitle{\insertshortsubtitle: \insertsection}
		\framesubtitle{\insertsubsection: Vanilla AE vs. VAE}
		
		\begin{itemize}
			\item {\scriptsize\url{https://github.com/projeduc/ESI_ML/blob/main/demos/NN/TF_Autoencoder.ipynb}}
		\end{itemize}
		
		\begin{minipage}{0.47\textwidth} 
			\begin{center}
				Standard Autoencoder
			\end{center}\vskip-6pt
			\hgraphpage[\textwidth]{auto-enc.png}
		\end{minipage}
		%
		\begin{minipage}{0.47\textwidth}
			\begin{center}
				Variational Autoencoder
			\end{center}\vskip-6pt
			\hgraphpage[\textwidth]{auto-enc-var.png}
		\end{minipage}
		
	\end{frame}
	
	\begin{frame}
		\frametitle{\insertshortsubtitle: \insertsection}
		\framesubtitle{\insertsubsection: Architecture}
		
		\hgraphpage[\textwidth]{AE-var.pdf}
		
	\end{frame}
	
	
	\begin{frame}
		\frametitle{\insertshortsubtitle: \insertsection}
		\framesubtitle{\insertsubsection: Latent distribution}
		
		\begin{itemize}
			\item Instead of learning a single latent vector, the encoder learns:
			
			\[
			\mu(x), \sigma(x)
			\]
			
			\item The latent vector is sampled from:
			
			\[
			z \sim \mathcal{N}(\mu, \sigma^2)
			\]
			
			\item Reparameterization trick:
			
			\[
			z = \mu + \sigma \odot \epsilon
			\]
			
			\[
			\epsilon \sim \mathcal{N}(0, I)
			\]
			
			\item This allows gradient backpropagation through stochastic sampling.
			
		\end{itemize}
		
	\end{frame}
	
	
	\begin{frame}
		\frametitle{\insertshortsubtitle: \insertsection}
		\framesubtitle{\insertsubsection: Loss function}
		
		\begin{itemize}
			\item The objective contains:
			
			\begin{itemize}
				\item reconstruction loss,
				\item distribution regularization.
			\end{itemize}
			
			\item Final objective:
			
%			\[
%			J_{VAE}
%			=
%			J(x, \hat{x})
%			+
%			KL\left(
%			\mathcal{N}(\mu, \sigma^2)
%			\;||\;
%			\mathcal{N}(0, I)
%			\right)
%			\]
			
			\item KL divergence encourages the latent space to remain close to a standard normal distribution.
			
			\item This produces smoother latent representations and improves sampling.
			
		\end{itemize}
		
	\end{frame}
	
	
	\begin{frame}
		\frametitle{\insertshortsubtitle: \insertsection}
		\framesubtitle{\insertsubsection: Practical remarks}
		
		\begin{itemize}
			\item Without regularization, the learned distributions may become too broad.
			
			\item KL divergence constrains the latent distributions.
			
			\item The variance must remain positive:
			
			\[
			\sigma > 0
			\]
			
			\item In practice, the network learns:
			
			\[
			\log \sigma^2
			\]
			
			instead of directly learning $\sigma$.
			
		\end{itemize}
		
	\end{frame}
	
	\section{Deep generative models}
	
	\begin{frame}
		\frametitle{\insertshortsubtitle}
		\framesubtitle{\insertsection}
		
	\end{frame}
	
	\section{Generative Adversarial Networks (GANs)}
	
	\begin{frame}
		\frametitle{\insertshortsubtitle}
		\framesubtitle{\insertsection: Motivation}
		
		\begin{itemize}
			\item Generative models learn the underlying distribution of the data.
			\item Their goal is to generate new realistic samples.
			\item Applications:
			\begin{itemize}
				\item Image generation
				\item Face synthesis
				\item Super-resolution
				\item Style transfer
				\item Data augmentation
			\end{itemize}
			\item GANs were introduced by \cite{2014-goodfellow-al}.
		\end{itemize}
		
	\end{frame}
	
	\begin{frame}
		\frametitle{\insertshortsubtitle}
		\framesubtitle{\insertsection: Main idea}
		
		\begin{block}{Adversarial learning}
			Two neural networks compete against each other:
		\end{block}
		
		\begin{itemize}
			\item \textbf{Generator (G)}
			\begin{itemize}
				\item Generates fake samples from random noise.
				\item Tries to fool the discriminator.
			\end{itemize}
			
			\item \textbf{Discriminator (D)}
			\begin{itemize}
				\item Receives real and generated samples.
				\item Predicts whether a sample is real or fake.
			\end{itemize}
		\end{itemize}
		
		\begin{center}
			\textbf{Generator vs Discriminator}
		\end{center}
		
	\end{frame}
	
	\begin{frame}
		\frametitle{\insertshortsubtitle}
		\framesubtitle{\insertsection: Architecture}
		
		\hgraphpage[0.95\textwidth]{GAN.pdf}
		
	\end{frame}
	
	\begin{frame}
		\frametitle{\insertshortsubtitle}
		\framesubtitle{\insertsection: Generator}
		
		\begin{itemize}
			\item Input:
			\begin{itemize}
				\item Random vector \(z\)
				\item Usually sampled from a normal or uniform distribution
			\end{itemize}
			
			\item Output:
			\begin{itemize}
				\item Synthetic sample \(\hat{x}\)
			\end{itemize}
			
			\item Goal:
			\begin{itemize}
				\item Produce realistic data
				\item Fool the discriminator
			\end{itemize}
		\end{itemize}
		
		\begin{align*}
			z \sim \mathcal{N}(0,1)
		\end{align*}
		
	\end{frame}
	
	\begin{frame}
		\frametitle{\insertshortsubtitle}
		\framesubtitle{\insertsection: Discriminator}
		
		\begin{itemize}
			\item Binary classifier:
			\begin{itemize}
				\item Real sample \(\rightarrow 1\)
				\item Fake sample \(\rightarrow 0\)
			\end{itemize}
			
			\item Input:
			\begin{itemize}
				\item Real data \(x\)
				\item Generated data \(G(z)\)
			\end{itemize}
			
			\item Output:
			\begin{itemize}
				\item Probability that the sample is real
			\end{itemize}
		\end{itemize}
		
		\begin{align*}
			D(x) \in [0,1]
		\end{align*}
		
	\end{frame}
	
	\begin{frame}
		\frametitle{\insertshortsubtitle}
		\framesubtitle{\insertsection: Objective function}
		
		\begin{itemize}
			\item The discriminator tries to maximize the classification accuracy.
			\item The generator tries to minimize the discriminator's success.
			\item Training is formulated as a minimax game.
		\end{itemize}
		
		\begin{align*}
			\min_G \max_D V(D,G) =
			& \mathbb{E}_{x \sim p_{data}(x)}[\log D(x)] \\
			+ &
			\mathbb{E}_{z \sim p_z(z)}[\log(1 - D(G(z)))]
		\end{align*}
		
	\end{frame}
	
	\begin{frame}
		\frametitle{\insertshortsubtitle}
		\framesubtitle{\insertsection: Training process}
		
		\begin{enumerate}
			\item Sample real data from the dataset.
			
			\item Generate fake samples using the generator.
			
			\item Train the discriminator:
			\begin{itemize}
				\item Real samples labeled as 1
				\item Fake samples labeled as 0
			\end{itemize}
			
			\item Train the generator:
			\begin{itemize}
				\item Update \(G\) to fool \(D\)
			\end{itemize}
			
			\item Repeat until convergence.
		\end{enumerate}
		
	\end{frame}
	
	\begin{frame}
		\frametitle{\insertshortsubtitle}
		\framesubtitle{\insertsection: Challenges}
		
		\begin{itemize}
			\item Training instability
			\item Difficult convergence
			\item Sensitive hyperparameters
			\item Vanishing gradients
			\item \textbf{Mode collapse}
			\begin{itemize}
				\item The generator produces limited varieties of samples.
			\end{itemize}
		\end{itemize}
		
	\end{frame}
	
	\begin{frame}
		\frametitle{\insertshortsubtitle}
		\framesubtitle{\insertsection: Variants}
		
		\begin{itemize}
			\item \textbf{DCGAN}
			\begin{itemize}
				\item Uses convolutional neural networks.
			\end{itemize}
			
			\item \textbf{Conditional GAN (CGAN)}
			\begin{itemize}
				\item Generation conditioned on labels or text.
			\end{itemize}
			
			\item \textbf{CycleGAN}
			\begin{itemize}
				\item Image-to-image translation.
			\end{itemize}
			
			\item \textbf{StyleGAN}
			\begin{itemize}
				\item High-quality face generation.
			\end{itemize}
		\end{itemize}
		
	\end{frame}
	
	\begin{frame}
		\frametitle{\insertshortsubtitle}
		\framesubtitle{\insertsection: Applications}
		
		\begin{itemize}
			\item Face generation
			\item Deepfake synthesis
			\item Image super-resolution
			\item Art generation
			\item Medical image synthesis
			\item Data augmentation
			\item Image restoration
		\end{itemize}
		
	\end{frame}
	
%	\begin{frame}
%		\frametitle{\insertshortsubtitle}
%		\framesubtitle{\insertsection: Example}
%		
%		\begin{center}
%			\hgraphpage[0.9\textwidth]{GAN-example.jpg}
%		\end{center}
%		
%	\end{frame}
	
	
	
	
	\insertbibliography{DL05L-Unsupervised}{*}
	
\end{document}